% Generated draft: requires mathematical and source review.
\documentclass[11pt]{article}
\usepackage[margin=25mm]{geometry}
\usepackage{fontspec}
\setmainfont{Latin Modern Roman}
\usepackage{amsmath,amssymb,mathtools,unicode-math}
\setmathfont{Latin Modern Math}
\usepackage{enumitem,needspace,xurl,hyperref}
\setlength{\parindent}{0pt}
\setlength{\parskip}{5pt}
\setlength{\emergencystretch}{3em}
\newcommand{\N}{\mathbb{N}}
\newcommand{\Z}{\mathbb{Z}}
\newcommand{\Q}{\mathbb{Q}}
\newcommand{\R}{\mathbb{R}}
\newcommand{\C}{\mathbb{C}}
\newcommand{\F}{\mathbb{F}}
\newcommand{\A}{\mathbb{A}}
\newcommand{\PP}{\mathbb{P}}
\newcommand{\E}{\mathbb{E}}
\newcommand{\eps}{\varepsilon}
\newcommand{\dd}{\,\mathrm{d}}
\begin{document}
\section*{UnsolvedMath ID: 30000091}
\textbf{Source problem code:} catalogue entry
\subsection{Definitions and mathematical statement}
\textbf{Source mismatch: no verified problem is promoted under this citation.} The retained original asks whether an impenetrable scatterer is uniquely determined by the far-field pattern for a single incident plane wave. The supplied verification says that the attached combinatorial-commutative-algebra report contains no such inverse-scattering problem and that the statement appears shifted from the following report. It marks the entry unrecoverable. No replacement report identity is supplied, so none is guessed.

The retained question can be explained only conditionally. In a scalar acoustic model, let $D\subset\R^m$, $m=2$ or $3$, be a bounded obstacle, let $k>0$ be a wavenumber and $d\in S^{m-1}$ an incident direction, and set
\[
 u^{\mathrm{inc}}(x)=e^{\mathrm i kx\cdot d},\qquad
 (\Delta+k^2)(u^{\mathrm{inc}}+u^{\mathrm{sc}})=0
 \quad\text{outside }\overline D.
\]
A chosen obstacle boundary condition and the outgoing radiation condition determine a far-field amplitude $u_D^\infty(\widehat x;k,d)$ as a function of observation direction. For a fixed admissible obstacle class and fixed boundary-condition convention, the single-wave uniqueness target would be
\[
 u_{D_1}^\infty(\widehat x;k,d)=u_{D_2}^\infty(\widehat x;k,d)
 \quad\forall\widehat x\in S^{m-1}
 \quad\Longrightarrow\quad D_1=D_2,
\]
with one fixed $(k,d)$. The original does not specify acoustic versus electromagnetic scattering, dimension, obstacle regularity, connectedness of the exterior, boundary condition, fixed-frequency quantification, or complex versus phaseless data. Thus this display is an illustrative interpretation of the retained text, not a verified formulation attributable to the attached report.

The supplied literature triage, checked 21 August 2026, instead discusses toric syzygies: strong positivity for sufficiently high tensor powers, with a borderline question for arbitrary ample line bundles not settled by those bounds. That assessment and its references are unrelated to the retained scattering statement. Neither its open label nor its claimed findings provide a literature status for single-wave inverse scattering.
\subsection{Short English statement}
The original single-plane-wave obstacle-uniqueness question is retained for provenance, but its report citation and literature assessment concern commutative algebra. The record is explicitly unrecoverable, so no verified scattering conjecture or literature status is assigned under the attached source.
\subsection{Sources}
\begin{itemize}
\item Local UnsolvedMath registry, permanent ID 30000091, exact source code catalogue entry; \emph{Single-Wave Uniqueness for Impenetrable Scatterers}. Proposed year 2004, registry status open, statement\_status unrecoverable. Supplied verification explicitly reports a shifted-statement/source mismatch and promotes no clean statement.
\item Attached, mismatched report citation: \emph{Combinatorial Commutative Algebra (2005)}, Oberwolfach Reports, DOI 10.4171/owr/2004/32, pp. 1703--1746, \url{https://doi.org/10.4171/owr/2004/32}. It is not supplied as a reliable inverse-scattering source; title year differs from DOI/proposed year.
\item Supplied but unrelated literature assessment checked 2026-08-21: ``Schenck--Smith, Syzygies of projective toric varieties'', \url{https://arxiv.org/abs/math/0308205}; ``Hering--Schenck--Smith, multigraded regularity'', \url{https://arxiv.org/abs/math/0502240}; ``Original Oberwolfach report'', \url{https://doi.org/10.4171/owr/2004/32}. These concern toric syzygies and ample bundles, not the scattering question. The dated OWR-derived curation requests independent verification.
\item No authoritative inverse-scattering citation or exact identity of the purported following report is supplied.
\end{itemize}
\end{document}
